% =============================================================================
% 论文类型定位与修订路径分析：从教师批评到学术重构
% =============================================================================

\documentclass[UTF8,a4paper,12pt,fontset=fandol]{ctexart}

\usepackage{geometry}
\geometry{left=2.2cm,right=2.2cm,top=2.2cm,bottom=2.2cm}
\usepackage{graphicx}
\usepackage{booktabs}
\usepackage{hyperref}
\usepackage{amsmath,amssymb}
\usepackage{setspace}
\setstretch{1.3}
% References handled via thebibliography

\title{\heiti\zihao{2} 学术论文的类型定位与修订路径分析 \\
       \large ——以"信用卡欺诈检测深度学习研究"课程作业为例}
\date{}

\begin{document}

\maketitle

\begin{abstract}
\noindent
课程大作业中的学术论文常常面临一个根本性的定位困境：一篇论文能否同时承担"应用验证"和"算法创新"两个任务？本文以一个具体的课程大作业案例——《深度神经网络在信用卡欺诈检测中的应用研究》——为分析对象，解剖其收到的四条教师批评背后的深层学术逻辑。分析表明，四条批评共同指向一个核心问题：论文的类型定位模糊导致其既不能满足应用研究对"外部有效性"（external validity）的要求，也无法达到算法研究对"内部机制"（internal mechanism）的探究深度。在此基础上，本文系统比较了三条可行的修订路径（聚焦FT-Transformer注意力塌缩机制、换文本数据集、聚焦DeepFM与树同构性），并从学术贡献增量、工作量可行性、叙事统一性三个维度论证了"聚焦FT-Transformer注意力塌缩机制"为最优路径。本文的讨论虽然基于一个具体案例，但其关于学术论文类型边界、归纳偏置相容性、以及创新层级的分析框架对更广泛的课程论文选题与修订具有参考意义。
\end{abstract}

% =============================================================================
\section{问题的提出}

一篇课程大作业论文收到老师四条批评后，最常见的反应是逐条修补——加内容回应第一条，删内容回应第二条，调参数回应第三条。但四条批评之间可能存在一个共同的逻辑根源，若不找到这个根源，修补将是零散的、甚至相互矛盾的。

本文分析的原案例为一篇商业智能课程的大作业论文，标题为《深度神经网络在信用卡欺诈检测中的应用研究——基于树归纳偏置迁移的深度表格模型》，使用Kaggle 2023年信用卡欺诈检测数据集（29维PCA匿名化特征，0.5\%欺诈率），比较了XGBoost、FT-Transformer、DeepFM三个基座模型及三个蒸馏变体的性能，并在SHAP、注意力权重和置换重要性三个维度进行了可解释性分析。全文27页。

该论文收到老师如下四条批评：

\begin{enumerate}
    \item \textbf{讲得太杂}：6个模型横向对比，既像应用研究（欺诈检测）又像算法研究（架构比较），两边都不够深。
    \item \textbf{应用研究不像应用研究}：工业界信用卡欺诈检测主要使用树模型（XGBoost/LightGBM），深度神经网络在该领域没有实际商业产出，比较本身缺乏应用意义。
    \item \textbf{算法研究不够深}：Sparsemax替代Softmax和知识蒸馏都是现有技术的组合应用，缺乏对"为什么有效/为什么失效"的机制层面深入探究。
    \item \textbf{建议}：不必跟树模型比较，聚焦一个模型深入做，或换文本数据集。
\end{enumerate}

这四条批评分别指向论文的叙事结构、应用价值、创新深度和实验设计，但它们之间是否存在一条统一的逻辑线索？本文将论证：四条批评共同指向一个根本问题——\textbf{论文的类型定位（paper type）模糊}。应用研究、算法改进和机制探究是三种不同的学术写作类型，各自有其独立的论证逻辑、评估标准和读者期待。当一篇论文试图同时满足三种类型的标准时，它往往会三种都不满足。

% =============================================================================
\section{批评背后的深层学术逻辑}

\subsection{学术论文的类型边界}

在机器学习和数据科学领域，论文可被大致分为三种类型，各自拥有不同的论证结构、评估标准和贡献定义\cite{wagstaff2012machine}：

\begin{table}[htbp]
\centering
\caption{三类学术论文的比较框架}
\label{tab:paper_types}
\begin{tabular}{p{2.5cm}p{4cm}p{4cm}p{4cm}}
\toprule
\textbf{维度} & \textbf{应用型论文} & \textbf{算法型论文} & \textbf{机制型论文} \\
\midrule
核心问题 & 方法X能否解决业务问题Y？ & 算法A是否比算法B更好？ & 算法A为什么有效/失效？ \\
贡献定义 & 验证外部有效性 & 提出新算法/方法 & 发现新机制/规律 \\
关键证据 & 业务指标改善 & 多数据集基准提升 & 消融实验，机制解释 \\
读者期待 & 可操作的建议 & 可复现的提升 & 可迁移的洞察 \\
典型结构 & 问题→方法→结果→讨论 & 背景→算法→实验→分析 & 观察→假设→验证→理论 \\
\bottomrule
\end{tabular}
\end{table}

将原案例论文映射到此框架中，可以发现其核心矛盾：论文的标题（"……应用研究"）和数据集选择（信用卡欺诈检测）指向\textbf{应用型}定位；其实验设计（6模型横向对比、多指标评估）指向\textbf{算法型}定位；其Sparsemax替换和TreeReg蒸馏分析隐含了\textbf{机制型}贡献的意图。三种类型被压缩在同一篇论文中，导致每一类型的核心问题都没有得到充分回答。

\subsection{应用型论文的"外部有效性"门槛}

应用型论文的论证核心是\textbf{外部有效性}（external validity）：在实验室环境中观察到的方法优势，能否迁移到真实业务环境中\cite{sculley2015hidden}？

对信用卡欺诈检测而言，工业界的真实环境有以下特征：（1）数据是流式的，模型需要频繁更新；（2）特征具有明确的业务语义（交易金额、商户类型、时间戳），可供业务专家理解和审计；（3）决策需要可解释性以满足监管合规要求（如GDPR第22条）；（4）树模型（XGBoost/LightGBM）已经是成熟的生产基线，任何替代方案需要在至少一个关键维度（精度、速度、可解释性、维护成本）上提供显著优势。

原案例论文在这些维度上均未提供充分的论证。它使用了一个PCA匿名化的数据集（丧失了业务语义），在单次静态划分上报告了AUPRC差异，但未讨论推理延迟、模型更新成本或合规审计能力。更重要的是，论文未回答"如果一个风控团队已经在用XGBoost，他们为什么要换成DeepFM+KD"这一核心业务问题——因为0.0164的AUPRC提升（0.9495 vs 0.9331）在商业场景中无法单独构成迁移决策的充分理由。

这解释了批评（2）的深层逻辑：\textbf{不是"DNN在欺诈检测上没有价值"，而是"论文没有论证DNN在欺诈检测上的价值"。}要论证这一价值，需要的不是6个模型的AUPRC比较，而是对精度提升如何转化为业务收益（如减少欺诈损失金额）的量化估算。

\subsection{算法型论文的"内在有效性"门槛}

算法型论文的论证核心是\textbf{内在有效性}（internal validity）：观察到的性能提升是否确实归因于所提出的算法改进，而非其他混淆因素\cite{henderson2018deep}？

原案例论文在内在有效性上面临两个挑战。第一，6个模型并非在同一条件下比较：FT-Transformer基线使用$d_{model}=64$, $n_{layers}=3$，而KD变体使用$d_{model}=48$, $n_{layers}=2$——架构参数差异使性能差异无法被唯一归因于蒸馏策略。第二，单一随机种子（seed=42）下0.0164的AUPRC差异不具有统计显著性——无法排除"另一个随机初始化可能产生不同排名"的替代解释。

这补充了批评（3）的逻辑：\textbf{不仅创新深度不够，而且因果推断的严谨性也不够。}算法型论文不能仅靠"我们做了X，然后性能从Y提升到Z"来建立贡献——它还需要排除替代解释、进行统计检验、建立效应量（effect size）的置信区间。

\subsection{机制型论文的"解释深度"门槛}

机制型论文的核心贡献在于回答"为什么"（why）而非"什么"（what）或"是否"（whether）。这需要三个层次的论证：（1）\textbf{现象确认}——可复现地记录某种行为模式；（2）\textbf{机制假设}——提出可检验的因果假说；（3）\textbf{机制验证}——设计控制实验排除替代假设。

原案例论文在现象确认层面已经积累了重要证据——Softmax注意力塌缩、Sparsemax修复效果、TreeReg破坏效应——但它将这些证据分散在27页的多条叙事线中，未提炼为一个有因果结构的整体。具体而言，它缺少：（a）塌缩的根因判定（PCA空间 vs Focal Loss vs 架构超参），（b）修复效应的机制分解（Sparsemax稀疏化 vs 熵调节，KD校准路径 vs 特征重排路径），（c）TreeReg失效的强度-效应剂量关系。

这解释了批评（1）和（3）的共同指向：\textbf{论文不是素材不够，而是叙事结构未能将这些素材组织为一个"回答为什么"的因果链条。}

% =============================================================================
\section{三条修订路径的比较分析}

基于上述诊断，以下系统比较三条修订路径。比较框架包含四个维度：学术贡献增量、工作量要求、与现有代码/数据的兼容性、核心风险的缓解。

\subsection{路径A：聚焦FT-Transformer注意力塌缩机制（推荐）}

\textbf{核心策略}：将论文从"6模型横向对比"重构为"自注意力在PCA匿名化表格数据上的塌缩诊断、修复与结构性拮抗发现"。删除或降级DeepFM和XGBoost（XGBoost仅作为KD教师保留），将全部实验资源集中在四个FTT变体（Softmax/Sparsemax/+KD/+KD+TreeReg）的消融因果链上。

\textbf{新的核心问题}：当表格数据的特征丧失了语义锚点（如PCA匿名化），自注意力机制是否会发生系统性塌缩？若会，不同的修复策略（归一化函数替换、输出层校准注入、嵌入层先验注入）各自的效果边界和失败模式是什么？

\textbf{学术贡献增量}：

\begin{itemize}
    \item \textbf{诊断贡献}：提供自注意力塌缩的完整实证证据链（注意力权重分布、多头多样性、嵌入低秩、Q-K内积方差），并排除Focal Loss为混淆因素（BCE消融）。
    \item \textbf{方法论贡献}：通过A→B→C→D的消融设计，首次将Sparsemax、KD和TreeReg的独立效应精确解耦——KD不通过重塑注意力分布起作用，TreeReg对注意力的破坏是结构性替代而非调参问题。
    \item \textbf{理论贡献}：提出"结构性替代"（Structural Substitution）概念——当外部注入的归纳偏置与目标架构的内在归纳偏置不可通约时，前者可能是在替代后者的核心功能，而非良性引导。这一概念可迁移至更广泛的归纳偏置迁移研究。
\end{itemize}

\textbf{工作量评估}：约5--7天。论文文字需重写（现有素材70\%可复用），可能需要1--2个补充分析实验（BCE消融、Q-K内积分析、注意力熵对比），均在已训练模型上运行，无需重新训练主模型。

\textbf{风险与缓解}：
\begin{itemize}
    \item \textit{风险}：BCE消融可能显示Focal Loss是塌缩主因而非PCA空间——这将削弱"PCA空间→塌缩"的因果关系，但为论文增加一层更有趣的细微发现（"塌缩是PCA+FocalLoss的交互效应"）。
    \item \textit{缓解}：论文的"结构性替代"核心论点不依赖塌缩的具体原因——只要TreeReg与注意力之间存在机制层面的冲突，该理论贡献即成立。
\end{itemize}

\subsection{路径B：换文本数据集}

\textbf{核心策略}：放弃PCA匿名化的信用卡欺诈数据集，换用包含文本字段的数据集（如IEEE-CIS Fraud Detection或Yelp评论欺诈检测），利用预训练语言模型处理文本语义信号。

\textbf{优点}：深度学习（Transformer）在文本数据上天然优于树模型，应用价值立得住。论文可以自然地使用BERT/RoBERTa等预训练模型，方法选择具有内部一致性。

\textbf{缺点}：
\begin{itemize}
    \item \textbf{全部代码作废}：所有数据管道、模型实现、训练和评估代码需要从头编写。
    \item \textbf{工作量过大}：约需2--3周全时工作，对1周的修订窗口不可行。
    \item \textbf{丧失现有发现的学术价值}：Sparsemax修复、TreeReg拮抗等发现的诊断价值远超其在任何应用数据集上的性能数字——这些发现对表格深度学习的理论讨论具有持久的参考意义。
\end{itemize}

\subsection{路径C：聚焦DeepFM与树同构性}

\textbf{核心策略}：以DeepFM为唯一研究对象，论证其FM一阶线性项$w_i x_i$与树模型的轴对齐分裂$x_i > t$之间的数学同构性，并通过蒸馏实验展示如何利用这一同构性实现高效的归纳偏置迁移。

\textbf{优点}：DeepFM是原论文中AUPRC最高的模型（0.9495），且FM-树同构性在理论上具有真正的学术深度。

\textbf{缺点}：
\begin{itemize}
    \item \textbf{缺乏"负结果"的叙事张力}：DeepFM+KD成功、DeepFM基线已接近XGBoost——论文将变成"我们的方法有效"的线性叙事，缺乏原论文中TreeReg破坏带来的戏剧性张力和深层洞察。
    \item \textbf{FM-树同构性可能被视为"显然"的}：审稿人可能指出"既然FM一阶项就是线性回归，线性回归和轴对齐分裂的同构性在数学上是显然的"——这将削弱论文的novelty声称。
    \item \textbf{与路径A相比，理论贡献更窄}：FM-树同构性是一个关于FM架构的特定发现，而Sparesemax+KD+TreeReg的三角张力揭示了更一般的关于归纳偏置相容性的问题。
\end{itemize}

\subsection{综合比较}

\begin{table}[htbp]
\centering
\caption{三条修订路径的多维度比较}
\label{tab:path_comparison}
\begin{tabular}{lccc}
\toprule
\textbf{维度} & \textbf{路径A（FTT塌缩）} & \textbf{路径B（文本数据）} & \textbf{路径C（DeepFM）} \\
\midrule
学术贡献增量 & 高（三层贡献） & 中（应用创新） & 中（架构特定） \\
理论贡献可迁移性 & 高（归纳偏置相容性） & 低 & 中（FM-树同构） \\
工作量（1周） & 可行 & 不可行（3周+） & 可行 \\
代码复用率 & 70--80\% & 0\% & 60--70\% \\
叙事张力 & 高（成功→失败→洞察） & 低（线性叙事） & 中 \\
风险 & 中（BCE消融可能弱化塌缩归因） & 高（全部从零开始） & 中（同构性可能被视为显然） \\
\bottomrule
\end{tabular}
\end{table}

% =============================================================================
\section{路径A的详细论证}

\subsection{为什么"注意力塌缩"是一个好的核心问题}

一个好的学术核心问题需要同时满足四个条件：（1）\textbf{清晰性}——可以被精确表述和操作化；（2）\textbf{未知性}——现有文献中缺乏充分回答；（3）\textbf{重要性}——答案对领域的理论或方法讨论有意义；（4）\textbf{可答性}——在现有数据和方法约束下可以给出有说服力的答案。

"自注意力在PCA匿名化特征空间上为什么会塌缩？不同的修复策略为什么有截然不同的效果边界？"这一问题满足以上四个条件：

\textbf{清晰性}：塌缩有严格的形式化定义（注意力权重与均匀分布的距离$\epsilon \ll 1/d$），可操作化为JS散度、标准差、多头余弦相似度等多个指标。

\textbf{未知性}：现有文献中，Gorishniy等人\cite{gorishniy2021revisiting}和Grinsztajn等人\cite{grinsztajn2022tree}报告了FT-Transformer在表格数据上的整体表现，但未专门分析其注意力权重的分布特性。Martins和Astudillo\cite{martins2016sparsemax}提出了Sparsemax的数学框架，但未在表格数据上进行消融对照实验。关于TreeReg与注意力机制之间的结构性拮抗，目前文献中无相关报道。

\textbf{重要性}：自注意力在NLP和CV中的范式性成功已导致其被不加检验地移植到各种数据模态上。如果某些类型的数据（如PCA匿名化的表格数据）在前提条件上就不满足自注意力的运行假设，那么该发现对表格深度学习的实践选型具有警示价值。此外，"结构性替代"的发现（TreeReg的失败模式）触及了一个更一般的问题：在迁移学习/蒸馏中，我们如何判断所迁移的知识形式（form）与接收架构的归纳偏置是兼容的？

\textbf{可答性}：四个FTT变体（Softmax/Sparsemax/+KD/+KD+TreeReg）的训练已完成，消融因果链的因果推断有效性已经建立。现有数据和方法可以支撑一个有说服力的答案。

\subsection{叙事结构的重构逻辑}

新论文的三幕叙事结构设计如下：

\begin{description}
    \item[第一幕：诊断] 使用Softmax的FT-Transformer发生了什么？定量证据→机制解释→根因判定。引导读者从"观察到奇怪的现象"走向"理解现象发生的原因"。
    \item[第二幕：修复] 三种不同侵入程度的修复策略被依次施加——Sparsemax（仅替换归一化，最低侵入）→KD（输出层间接注入，中等侵入）→TreeReg（嵌入层直接约束，最高侵入）。每种修复的策略逻辑被明确阐述，其成败不再是一个"惊喜"，而是策略逻辑的被检验后果。
    \item[第三幕：洞察] 三种修复的差异化结果在"注意力熵-性能"平面上的拓扑关系被分析。Sparsemax和KD的成功路径（"在不改变注意力行为的前提下改善输出质量"）与TreeReg的失败路径（"在改变注意力行为的同时破坏输出质量"）的对比，引出"结构性替代"框架。
\end{description}

这一叙事结构的关键优势在于：\textbf{它将原论文中分散在6个模型和27页中的观察，集中组织为一个有因果方向的因果链}。读者不需要同时记忆6个模型的性能数字——他们只需要追踪4个变体从A到D的演变轨迹。

\subsection{删除内容的逻辑}

修订方案建议删除以下内容（或降级为附录）：

\begin{itemize}
    \item \textbf{DeepFM完整架构描述}：DeepFM在论文中不再是"比较对象"，而是"跨架构参考"。其FM-树同构性可在讨论章节中以一个段落的篇幅提及，作为"注意力型架构的结构性替代问题在交互型架构上可能不存在"的对照论证。
    \item \textbf{TabNet}：在原论文中已占篇幅较少，且其Sequential Attention与FTT的全局自注意力的关系不清晰。删除不引入任何叙事损失。
    \item \textbf{"模型选择建议"表}：这是一个应用型论文的标配元素——告诉实践者"什么场景选什么模型"。但在机制型论文中，它是不合适的——因为论文的核心贡献不是指导实践者选模型，而是告诉研究者"自注意力在什么条件下失效以及为什么"。
    \item \textbf{SMOTE/ClassWeight/PlainCE消融}：这些是不平衡处理的消融——它们在应用型或算法型论文中是相关的，但在机制型论文中偏离核心叙事。
\end{itemize}

\subsection{核心风险与应对}

方案A的主要风险是一个经验性问题：如果BCE消融显示Focal Loss是导致注意力塌缩的主要原因（而非PCA空间），则"PCA空间→塌缩"的因果归因将需要修正。

应对策略有两种。\textbf{策略一（接受复杂性）}：若BCE的确产生了更分散的注意力分布，论文的诊断将从"PCA空间导致塌缩"修正为"PCA空间+Focal Loss交互效应导致塌缩"——这实际上是一个更细致、更符合实际情况的结论。\textbf{策略二（防御性重述）}：论文可在BCE消融实验之前，将塌缩的因果归因从断言式（"PCA空间导致塌缩"）调整为假设式（"PCA空间是塌缩的可能原因，BCE消融将检验Focal Loss是否为混淆因素"）。这两种策略都不损害论文的核心贡献——"结构性替代"的发现完全独立于塌缩的具体原因。

% =============================================================================
\section{预期的学术贡献}

若成功按照路径A重构论文，预期将产生以下三个层级的学术贡献：

\subsection{第一层：发现级贡献——注意力塌缩的实证诊断}

在具体的实验条件下（PCA匿名化、29维特征、0.5\%欺诈率），本文提供了自注意力发生塌缩的完整实证证据。这一发现的学术价值不在于"这是第一次有人发现注意力会塌缩"（塌缩的数学可能性在Sparsemax原文中已被指出），而在于"这是第一次有人系统地记录了塌缩的实证特征并排除了主要的替代解释"。这为后续研究提供了一个可参照的诊断框架——其他研究者可以问："在我的数据上，注意力分布的标准差是否也接近零？如果是，我的嵌入层的SVD谱是否也高度集中？"这种可操作化的观察框架的经验积累，是领域从"模糊直觉"走向"精确理解"的关键路径。

\subsection{第二层：方法论贡献——消融因果链的设计}

A→B→C→D的累加式消融设计，精确解耦了Sparsemax（归一化函数替换）、KD（输出层先验注入）和TreeReg（嵌入层先验注入）的独立效应。特别是C→D的对比——"D在C的基础上唯一添加了TreeReg"——使得-3.73pp的性能下降可以被唯一归因于TreeReg。这种"一个变体只改变一个变量"的消融哲学本身是常规的，但本文将其应用于蒸馏策略的解耦，产生了一个方法论示范：如何系统地"拆解"一个由多个组件构成的蒸馏方案，以确定每个组件的独立贡献方向。

\subsection{第三层：理论贡献——结构性替代框架}

"结构性替代"（Structural Substitution）概念的提出——当外部归纳偏置与目标架构的内在归纳偏置基于不可通约的概念基础时，强制性迁移前者可能导致后者的结构性崩溃——是一个可推广至超越本文具体实验情境的理论贡献。这一框架对以下研究问题具有潜在的解释力：为什么某些跨模态的知识迁移会失败？为什么某些多任务学习方法会导致性能退化而非提升？在什么条件下"注入先验"成为"限制能力"？

% =============================================================================
\section{结论}

本文剖析了四条课程作业批评背后的深层学术逻辑，论证了它们共同指向一个称为"论文类型定位模糊"的结构性问题。通过对三种论文类型（应用型、算法型、机制型）的核心论证逻辑的解析，本文展示了混用这些逻辑如何导致每一类型的要求都无法被满足。

在三条修订路径的比较中，本文论证了"聚焦FT-Transformer注意力塌缩机制"为最优路径：它保留了原论文中最有价值的诊断证据（Sparsemax修复、TreeReg拮抗），将它们重新组织为一条有因果方向的叙事链，并通过提出"结构性替代"概念将其提升为可推广的理论贡献。该路径在1周的修订窗口中可行，代码复用率70--80\%，且核心风险（BCE消融可能弱化塌缩因果归因）不会威胁论文的核心论点。

更一般地，本文希望传达的一个信息是：\textbf{课程论文的选题和修订不应该从"我能做什么实验"出发，而应该从"我想回答什么问题"出发。}当核心问题清晰时，实验设计的必要性、充分性和因果推断的有效性都将自然呈现。当核心问题模糊时，即使有6个模型、27页篇幅和0.9495的AUPRC，论文仍然可能缺乏学术说服力——因为读者不知道该被说服什么。

\begin{thebibliography}{99}

\bibitem{wagstaff2012machine}
K.~L.~Wagstaff.
\newblock {Machine Learning that Matters}.
\newblock In \emph{Proc. ICML}, 2012, pp.~529--536.

\bibitem{sculley2015hidden}
D.~Sculley, G.~Holt, D.~Golovin, E.~Davydov, T.~Phillips, D.~Ebner, V.~Chaudhary, M.~Young, J.-F.~Crespo, and D.~Dennison.
\newblock {Hidden Technical Debt in Machine Learning Systems}.
\newblock In \emph{Advances in NIPS}, 2015, pp.~2503--2511.

\bibitem{henderson2018deep}
P.~Henderson, R.~Islam, P.~Bachman, J.~Pineau, D.~Precup, and D.~Meger.
\newblock {Deep Reinforcement Learning that Matters}.
\newblock In \emph{Proc. AAAI}, 2018, pp.~3207--3214.

\bibitem{gorishniy2021revisiting}
Y.~Gorishniy, I.~Rubachev, V.~Khrulkov, and A.~Babenko.
\newblock {Revisiting Deep Learning Models for Tabular Data}.
\newblock In \emph{Advances in NIPS}, Vol.~34, 2021, pp.~18932--18943.

\bibitem{grinsztajn2022tree}
L.~Grinsztajn, E.~Oyallon, and G.~Varoquaux.
\newblock {Why do tree-based models still outperform deep learning on typical tabular data?}
\newblock In \emph{Advances in NIPS}, Vol.~35, 2022, pp.~507--520.

\bibitem{martins2016sparsemax}
A.~F.~T.~Martins and R.~F.~Astudillo.
\newblock {From Softmax to Sparsemax: A Sparse Model of Attention and Multi-Label Classification}.
\newblock In \emph{Proc. ICML}, 2016, pp.~1614--1623.

\end{thebibliography}

\end{document}
